\section{Deployment Guide}
\subsection{Raspberry Pi 4 Setup}
\label{subsec:rpi4_setup}

This subsection describes the hardware platform, software stack, and deployment
procedure used to validate the proposed AUTOSAR SecOC gateway with post-quantum
cryptography on real embedded hardware.

\paragraph{Hardware Platform}
\mbox{}\\
The target hardware platform is a Raspberry Pi~4 Model~B equipped with 4~GB of
RAM. The board provides sufficient computational resources to evaluate the
performance and feasibility of post-quantum cryptographic primitives in an
automotive gateway context. An Ethernet interface is used as the primary
communication channel for secured PDU transmission. A 32~GB microSD card is used
for the operating system and development environment. An external MCP2515 CAN
controller is optionally supported for future CAN bus integration, although it
is not required for the Ethernet-based validation presented in this work.

\begin{itemize}
	\item Raspberry Pi~4 (4~GB RAM)
	\item Ethernet network interface
	\item microSD card (32~GB)
	\item MCP2515 CAN controller (optional, future work)
\end{itemize}

\paragraph{Software Stack}
\mbox{}\\
Figure~\ref{fig:rpi_software_stack} illustrates the software stack deployed on
the Raspberry Pi. The AUTOSAR SecOC application, extended with post-quantum
cryptography, executes in user space and links against the Open Quantum Safe
\texttt{liboqs} library compiled for ARM64. The application runs on top of a
64-bit Linux kernel (version~6.1), provided by Raspberry Pi OS based on Debian
Bookworm.

\begin{figure}[H]
	\centering
	\includegraphics[width=0.4\linewidth]{images/Chapter3/pi_soft_stack}
	\caption{Raspberry Pi Software Stack}
	\label{fig:rpi_software_stack}
\end{figure}

This layered architecture enables a clear separation between the AUTOSAR
security logic, the post-quantum cryptographic primitives, and the underlying
operating system, closely reflecting a realistic automotive gateway deployment
model.

\paragraph{Build Environment Setup}
\mbox{}\\
Before building the project, the following dependencies must be installed on the
target platform.

\footnotesize
\begin{longtable}{
		>{\raggedright\arraybackslash}p{3.8cm}
		>{\raggedright\arraybackslash}p{5.2cm}
		>{\raggedright\arraybackslash}p{5.4cm}
	}
	\caption{Build Dependencies}
	\label{tab:build_dependencies} \\
	\hline
	\textbf{Dependency} &
	\textbf{Windows (MinGW)} &
	\textbf{Raspberry Pi (Linux)} \\
	\hline
	\endfirsthead

	\multicolumn{3}{c}{\textit{(Continued from previous page)}} \\ \hline
	\textbf{Dependency} &
	\textbf{Windows (MinGW)} &
	\textbf{Raspberry Pi (Linux)} \\
	\hline
	\endhead

	\hline
	\multicolumn{3}{r}{\textit{(Continued on next page)}} \\
	\endfoot

	\hline
	\endlastfoot

	GCC Compiler &
	MinGW-w64 (GCC 13+) &
	\texttt{sudo apt install build-essential} \\

	CMake &
	CMake 3.20+ &
	\texttt{sudo apt install cmake} \\

	Ninja Build &
	Ninja 1.10+ &
	\texttt{sudo apt install ninja-build} \\

	OpenSSL &
	OpenSSL 3.0+ (bundled) &
	\texttt{sudo apt install libssl-dev} \\

	Git &
	Git for Windows &
	\texttt{sudo apt install git} \\

\end{longtable}
\normalsize

\paragraph{Building liboqs Library}
\mbox{}\\
The Open Quantum Safe (liboqs) library must be compiled before building the main
project. The build script automatically detects the platform and applies
appropriate compiler flags.

\begin{verbatim}
# Navigate to liboqs directory
cd Autosar_SecOC/external/liboqs
mkdir -p build && cd build

# Raspberry Pi (ARM64) configuration
cmake -GNinja \
    -DCMAKE_BUILD_TYPE=Release \
    -DBUILD_SHARED_LIBS=OFF \
    -DOQS_USE_OPENSSL=ON \
    -DOQS_ENABLE_KEM_ml_kem_768=ON \
    -DOQS_ENABLE_SIG_ml_dsa_65=ON \
    -DCMAKE_C_FLAGS="-mcpu=cortex-a72 -O3 -fPIC" \
    ..

# Build the library
ninja
\end{verbatim}

The resulting static library (\texttt{liboqs.a}) is located in
\texttt{external/liboqs/build/lib/} and includes only the ML-KEM-768 and
ML-DSA-65 algorithms required for this implementation.

\paragraph{Building the Test Suite}
\mbox{}\\
The project uses Google Test (fetched automatically via CMake) for unit testing.
The following test executables are built:

\footnotesize
\begin{longtable}{
		>{\raggedright\arraybackslash}p{5.0cm}
		>{\raggedright\arraybackslash}p{9.4cm}
	}
	\caption{Test Executables}
	\label{tab:test_executables} \\
	\hline
	\textbf{Executable} &
	\textbf{Description} \\
	\hline
	\endfirsthead

	\multicolumn{2}{c}{\textit{(Continued from previous page)}} \\ \hline
	\textbf{Executable} &
	\textbf{Description} \\
	\hline
	\endhead

	\hline
	\multicolumn{2}{r}{\textit{(Continued on next page)}} \\
	\endfoot

	\hline
	\endlastfoot

	\texttt{test\_pqc\_standalone} &
	ML-KEM-768 and ML-DSA-65 standalone algorithm validation \\

	\texttt{test\_pqc\_secoc\_integration} &
	CSM layer integration with classical MAC comparison \\

	\texttt{test\_phase3\_complete} &
	Complete Ethernet gateway with ML-KEM + HKDF + ML-DSA \\

	\texttt{test\_phase4\_integration\_report} &
	AUTOSAR architecture validation report generator \\

	\texttt{AuthenticationTests} &
	MAC and signature generation tests \\

	\texttt{VerificationTests} &
	MAC and signature verification tests \\

	\texttt{FreshnessTests} &
	Freshness counter and replay protection tests \\

	\texttt{SecOCTests} &
	Transport Protocol layer callback tests \\

	\texttt{startOfReceptionTests} &
	Buffer management and reception initiation tests \\

	\texttt{PQC\_ComparisonTests} &
	Classical vs PQC performance comparison tests \\

\end{longtable}
\normalsize

\paragraph{Thesis Validation Script}
\mbox{}\\
The \texttt{build\_and\_run.sh} script automates the complete thesis validation
process. It executes four phases of testing and generates comprehensive reports.

\begin{verbatim}
# Run complete thesis validation
cd Autosar_SecOC
bash build_and_run.sh
\end{verbatim}

The script performs the following validation sequence:

\footnotesize
\begin{longtable}{
		>{\raggedright\arraybackslash}p{1.6cm}
		>{\raggedright\arraybackslash}p{4.8cm}
		>{\raggedright\arraybackslash}p{7.8cm}
	}
	\caption{Thesis Validation Phases}
	\label{tab:validation_phases} \\
	\hline
	\textbf{Phase} &
	\textbf{Name} &
	\textbf{Coverage} \\
	\hline
	\endfirsthead

	\multicolumn{3}{c}{\textit{(Continued from previous page)}} \\ \hline
	\textbf{Phase} &
	\textbf{Name} &
	\textbf{Coverage} \\
	\hline
	\endhead

	\hline
	\multicolumn{3}{r}{\textit{(Continued on next page)}} \\
	\endfoot

	\hline
	\endlastfoot

	1 &
	PQC Fundamentals &
	ML-KEM-768 KeyGen/Encaps/Decaps, ML-DSA-65 KeyGen/Sign/Verify \\

	2 &
	Google Test Suite &
	42 unit tests across 7 test suites (AuthenticationTests, VerificationTests, FreshnessTests, SecOCTests, startOfReceptionTests, PQC\_ComparisonTests, Phase3\_Complete\_Integration) \\

	3 &
	Ethernet Gateway &
	ML-KEM key exchange, HKDF session key derivation, ML-DSA signatures, full AUTOSAR stack (COM $\rightarrow$ PduR $\rightarrow$ SecOC $\rightarrow$ Csm $\rightarrow$ SoAd) \\

	4 &
	Architecture Report &
	AUTOSAR BSW layer architecture, TX/RX sequences, PDU structure visualization \\

\end{longtable}
\normalsize

\paragraph{Configuration Validation}
\mbox{}\\
The script automatically validates the PQC configuration before running tests.
The following parameters are checked:

\begin{verbatim}
# Required configuration in include/SecOC/SecOC_PQC_Cfg.h
SECOC_USE_PQC_MODE              TRUE    # Enable ML-DSA-65 signatures
SECOC_USE_MLKEM_KEY_EXCHANGE    TRUE    # Enable ML-KEM-768 key exchange
SECOC_ETHERNET_GATEWAY_MODE     TRUE    # Enable Ethernet gateway mode
SECOC_PQC_MAX_PDU_SIZE          8192    # Support 3,309-byte signatures
\end{verbatim}

\paragraph{Output Files}
\mbox{}\\
Upon successful completion, the validation script generates the following output
files:

\footnotesize
\begin{longtable}{
		>{\raggedright\arraybackslash}p{6.4cm}
		>{\raggedright\arraybackslash}p{8.0cm}
	}
	\caption{Validation Output Files}
	\label{tab:output_files} \\
	\hline
	\textbf{File} &
	\textbf{Description} \\
	\hline
	\endfirsthead

	\multicolumn{2}{c}{\textit{(Continued from previous page)}} \\ \hline
	\textbf{File} &
	\textbf{Description} \\
	\hline
	\endhead

	\hline
	\multicolumn{2}{r}{\textit{(Continued on next page)}} \\
	\endfoot

	\hline
	\endlastfoot

	\texttt{test\_summary.txt} &
	Comprehensive thesis validation report \\

	\texttt{test\_logs/phase1\_<timestamp>.txt} &
	PQC fundamentals test output \\

	\texttt{test\_logs/phase2\_<timestamp>.txt} &
	Google Test suite results \\

	\texttt{test\_logs/phase3\_<timestamp>.txt} &
	Ethernet gateway integration results \\

	\texttt{test\_logs/phase4\_<timestamp>.txt} &
	Architecture integration report \\

	\texttt{pqc\_standalone\_results.csv} &
	ML-KEM and ML-DSA performance measurements \\

	\texttt{pqc\_secoc\_integration\_results.csv} &
	SecOC integration performance data \\

\end{longtable}
\normalsize

\paragraph{Platform-Specific Build Commands}
\mbox{}\\
Table~\ref{tab:platform_build_commands} summarizes the build commands for each
supported platform.

\footnotesize
\begin{longtable}{
		>{\raggedright\arraybackslash}p{4.0cm}
		>{\raggedright\arraybackslash}p{5.0cm}
		>{\raggedright\arraybackslash}p{5.4cm}
	}
	\caption{Platform-Specific Build Commands}
	\label{tab:platform_build_commands} \\
	\hline
	\textbf{Operation} &
	\textbf{Windows (MinGW)} &
	\textbf{Raspberry Pi (Linux)} \\
	\hline
	\endfirsthead

	\multicolumn{3}{c}{\textit{(Continued from previous page)}} \\ \hline
	\textbf{Operation} &
	\textbf{Windows (MinGW)} &
	\textbf{Raspberry Pi (Linux)} \\
	\hline
	\endhead

	\hline
	\multicolumn{3}{r}{\textit{(Continued on next page)}} \\
	\endfoot

	\hline
	\endlastfoot

	CMake Generator &
	\texttt{-G "MinGW Makefiles"} &
	\texttt{-G "Unix Makefiles"} \\

	Build Command &
	\texttt{mingw32-make -j4} &
	\texttt{make -j4} \\

	Platform Flag &
	\texttt{-DWINDOWS} &
	\texttt{-DRASPBERRY\_PI -mcpu=cortex-a72} \\

	Socket Library &
	\texttt{-lws2\_32} &
	(POSIX sockets built-in) \\

	Ethernet Source &
	\texttt{ethernet\_windows.c} &
	\texttt{ethernet.c} \\

\end{longtable}
\normalsize

\paragraph{Raspberry Pi Deployment Procedure}
\mbox{}\\
The following steps deploy and validate the PQC Ethernet Gateway on a Raspberry
Pi~4:

\begin{enumerate}
	\item \textbf{Clone Repository:}
	\begin{verbatim}
git clone <repository-url>
cd computer_engineering_project/Autosar_SecOC
	\end{verbatim}

	\item \textbf{Install Dependencies:}
	\begin{verbatim}
sudo apt update
sudo apt install build-essential cmake ninja-build libssl-dev git
	\end{verbatim}

	\item \textbf{Run Validation:}
	\begin{verbatim}
bash build_and_run.sh
	\end{verbatim}

	\item \textbf{Verify Results:}
	\begin{verbatim}
cat test_summary.txt
	\end{verbatim}
\end{enumerate}

The validation script automatically detects the ARM64 platform and applies
Cortex-A72 optimizations during compilation.

